Este trabalho revisita o Trabalho de Graduação de W.~L.~S. de Paula
(ITA, 2003), que obteve seis estados de polarização para ondas
gravitacionais na teoria bimétrica de Visser com gráviton massivo, e
leva seus resultados até o observável de redes de temporização de
pulsares (PTAs). Refeita sem a aproximação de propagação quase nula, a
análise confirma exatamente as seis amplitudes de maré independentes.
Mostra também que a sexta é o traço da perturbação, um escalar fantasma
que permite ao modelo evitar a descontinuidade de van
Dam--Veltman--Zakharov. Na teoria consistente de Fierz--Pauli, restam
cinco polarizações, e os modos escalares transversal e longitudinal ficam
vinculados por $p_l=-(f_g/f)^2p_b$. As amplitudes de Newman--Penrose do
TG são contrações exatas do Riemann, mas sua leitura como amplitudes de
polarização e a classificação E(2) supõem propagação quase nula e só se
aplicam para $f\gg f_g$. No exemplo de baixa frequência do artigo de
2004, essa leitura erra cerca de 60\%. A supressão
$(f_g/f)^2$ dos modos adicionais encontrada no TG é exata para
amplitudes métricas comparáveis. Com os limites atuais sobre a massa,
porém, as assinaturas especificamente massivas só são apreciáveis na
banda de nanohertz. Mostra-se que a resposta de uma PTA depende da onda
apenas pela matriz de marés calculada no TG. Calculam-se as correlações
angulares dos setores tensorial e de helicidade zero, com termos dos
pulsares, e explica-se por simetria por que, no limiar $f=f_g$, todas as
polarizações produzem a mesma correlação por estado. Simulações
controladas mostram que redes pequenas aprendem pouco sobre a massa
(0,07 a 0,11 nat), que a média fixa das correlações sobre as frequências
elimina cerca de 95\% dessa informação e que a aproximação normal usual dos estimadores de
correlação multiplica por dez a taxa de falsas detecções de polarização
escalar.
